% group3_foundations_synthesis.tex
% GeoVac Group 3 Synthesis: Foundations / Spectral Graph Theory Arc
% Status: Project archive synthesis, May 2026
% Audience: mathematical physicists, spectral graph theorists, foundations community

\documentclass[aps,pra,twocolumn,superscriptaddress,longbibliography]{revtex4-2}

\usepackage{amsmath,amssymb,amsthm,graphicx,xcolor,bm,braket,booktabs}
\usepackage{hyperref}

\newtheorem{theorem}{Theorem}
\newtheorem{proposition}{Proposition}
\newtheorem{corollary}{Corollary}
\newtheorem{observation}{Observation}
\newtheorem{conjecture}{Conjecture}

\begin{document}

\title{Foundations of the Geometric Vacuum Framework:\\
A Synthesis of the Spectral Graph Theory Arc}
\thanks{Group~3 synthesis in the Geometric Vacuum series}

\author{J.~Loutey}
\affiliation{Independent Researcher, Kent, Washington}

\date{June 14, 2026}

%% ========================================================================
%% ABSTRACT
%% ========================================================================

\begin{abstract}
The Geometric Vacuum (GeoVac) framework is a discretization scheme for
non-relativistic quantum mechanics on natural geometries: the discrete
graph Laplacian on a packed lattice of quantum-number labels replaces
the continuous Laplace--Beltrami operator on the corresponding compact
manifold.  This synthesis paper assembles the framework's foundational
arc---the packing construction, the discrete
graph and its dynamics, the conformal equivalence with the unit $S^3$,
the exchange constant taxonomy, the angular sparsity theorem, the
Bargmann--Segal lattice, and the universal/Coulomb partition, together
with the \emph{reconvergence} node where these threads cohere (two-body
selection rules from the tensor-product triple, the period
classification, the Tannakian substrate and cosmic-Galois injection, and
the forced/free seam)---into a single coherent narrative.  The arc has a
dependency structure worth naming:\ a single taproot (the packing
construction, made physical by the Fock $S^3$ projection) fans out into
the framework's machinery, then reconverges at the periods, Tannakian,
and forced/free nodes, where the scattered structural findings cohere
into one statement about what the framework forces and what it consumes.  Four headline structural results frame the
program: (a) the discrete Fock graph is conformally equivalent to the
unit $S^3$ Laplace--Beltrami operator, verified by 18 independent
symbolic proofs; (b) the angular sparsity of two-body matrix elements
depends only on $l_{\max}$ and is bit-identical across five distinct
central potentials, qualifying as a potential-independent theorem
about spherical-fermion bases; (c) the three-dimensional isotropic
harmonic oscillator admits a discrete graph encoding on the Hardy
space $H^2(S^5)$ that is bit-exactly $\pi$-free in exact rational
arithmetic at every finite cutoff; (d) the framework's results split
cleanly along an A/D partition---an algebra-only universal sector that
transfers free across central potentials, and a Dirac-operator-specific
sector containing the spectrum, the calibration $\pi$, and the
candidate fine-structure-constant combination of Paper~2.  A structural reading we make explicit in this synthesis is that
the framework's substrate-level scale-freedom---the packing diameter
$d_{0}$ cancelling out of every spectral quantity, and physical
scales entering only through the Fock-projection focal length
$p_{0}$---is the geometric expression of the same principle that
conformal field theory~\cite{klebanov_pufu_safdi2011} embodies at
the physics level.  GeoVac provides a scale-free mathematical
substrate; CFT is a scale-free physics theory living on such a
substrate.  This is not a claim that GeoVac is a CFT; it is a
structural reflection that explains why the framework reproduces
continuum CFT-on-sphere observables bit-exactly (documented in the
operator-algebra arc, outside the foundations scope of this
synthesis).  We present
GeoVac as a discretization framework that exploits natural geometry,
not as a new foundation for quantum mechanics: the discrete graph and
the continuous Laplace--Beltrami operator are dual descriptions
connected by proven conformal maps, and the interpretive question of
which description is more fundamental is left open.  The framework's
concrete computational advantages---$O(V)$ sparsity, integer
eigenvalues on the unit $S^3$, angular momentum selection rules baked
into the basis, bit-exactly rational matrix elements on the
Bargmann--Segal lattice, and the universal sparsity bound---are
independent of that interpretive question.
\end{abstract}

\maketitle

%% ========================================================================
\section{Introduction}
\label{sec:intro}
%% ========================================================================

The standard formulation of non-relativistic quantum mechanics
presupposes a continuous configuration space and treats the
Schr\"odinger operator $-\tfrac{1}{2}\nabla^2 + V(r)$ as the
fundamental object.  For the hydrogen atom in particular, Fock showed
in 1935~\cite{fock1935} that the momentum-space Schr\"odinger equation
maps by stereographic projection onto the free Laplacian on the unit
three-sphere $S^3$.  This observation has been developed extensively
by Bargmann~\cite{bargmann1936}, Bander and Itzykson~\cite{bander_itzykson1966},
Avery~\cite{avery_book1989}, Aquilanti and
co-workers~\cite{aquilanti_caligiana2003}, and others working in the
hyperspherical-momentum-space tradition.  In each case, the $S^3$
structure is treated as an elegant reformulation of a known continuum
problem: a useful change of variables for analytic work, not a
foundation in its own right.

The Geometric Vacuum (GeoVac) framework
\cite{loutey_paper0,loutey_paper1,loutey_paper7} takes a complementary
computational stance.  Starting from a packing construction on an
isotropic sphere, it builds a finite directed graph whose nodes are
quantum-number labels $(n, l, m)$ and whose edges are transition
amplitudes generated by the $SU(2) \otimes SU(1,1)$ ladder algebra of
hydrogen.  The graph Laplacian on this lattice is dimensionless and
scale-invariant; physical energies emerge only upon imposing the Fock
energy-shell constraint $p_0^2 = -2E$, which acts as a stereographic
focal length.  The framework's central claim, established by Paper~7,
is that this discrete construction converges to the unit $S^3$
Laplace--Beltrami operator in the continuum limit.  The discrete and
continuous descriptions are connected by a proven conformal
equivalence, and we present them throughout as dual descriptions of
the same physics; the interpretive question of ontological priority is
explicitly left to the reader.

GeoVac differs from three nearby families of discretization schemes
in ways worth naming at the outset.  Lattice gauge theory discretizes
the spacetime continuum on a hypercubic grid; GeoVac discretizes the
quantum-number space and inherits the natural geometry of separable
quantum systems.  Tensor-network methods choose a basis adapted to
entanglement structure; GeoVac chooses a basis adapted to angular
momentum and the Fock conformal projection.  Algebraic quantum field
theory works at the level of operator algebras over Minkowski
spacetime; GeoVac is non-relativistic and operates on compact
Euclidean-signature manifolds.  The rhetoric rule that governs the
GeoVac papers (see the project's positioning statement) applies here:
GeoVac is not a new foundation for quantum mechanics, and the
synthesis presented in this paper does not assert ontological priority
of either the discrete or the continuum description.

The Group~3 (foundations) arc comprises eleven papers, organized by a
single dependency structure.  \emph{The taproot:} Paper~0
\cite{loutey_paper0} states the packing axiom and derives the angular
quantum numbers as packing indices; Paper~7~\cite{loutey_paper7} makes
it physical, proving the conformal equivalence with the unit $S^3$ via
18 independent symbolic verifications and extending to the
multi-electron $S^{3N-1}$ setting.  \emph{The fan-out} (the framework's
working machinery):\ Paper~1~\cite{loutey_paper1} sets up the spectral
graph theory and reproduces the Rydberg spectrum;
Paper~18~\cite{loutey_paper18} classifies transcendental content via a
six-tier exchange-constant taxonomy and the master Mellin engine;
Paper~22~\cite{loutey_paper22} states the potential-independent angular
sparsity theorem; Paper~24~\cite{loutey_paper24} constructs the
Bargmann--Segal lattice and the six-layer Coulomb/HO asymmetry;
Paper~31~\cite{loutey_paper31} states the $\mathcal{A}/D$
universal/Coulomb partition in spectral-triple language.  \emph{The
reconvergence} (where the branches loop back to a unified foundational
statement, \S\ref{sec:reconvergence}):\ Paper~54 derives two-body
selection rules from the tensor-product triple, Paper~55 classifies
every GeoVac period as cyclotomic mixed-Tate at level $\le 4$,
Paper~56 closes the cosmic-Galois injection $U^*_{\mathrm{GV}}
\hookrightarrow \mathcal{U}_4^{\mathrm{ab}} \rtimes SL_2$, and Paper~57
catalogues the forced/free seam.  An $O(V)$ time-evolution engine
(Paper~6~\cite{loutey_paper6}, now an archived standalone tool)
supplies the dynamical demonstration; this synthesis supersedes the
earlier single-document synthesis (Paper~21~\cite{loutey_paper21}).

Four headline structural results organize the rest of this synthesis.
\textbf{(H1)} The discrete Fock graph is conformally equivalent to the
unit $S^3$ at the level of the Laplace--Beltrami operator
(Paper~7~\cite{loutey_paper7}).  \textbf{(H2)} The angular sparsity
density of two-body matrix elements is potential-independent
(Paper~22~\cite{loutey_paper22}): a structural theorem of
spherical-fermion bases, not a Coulomb-specific accident.
\textbf{(H3)} The 3D harmonic oscillator admits a discrete graph
encoding on the Hardy sector of $S^5$ that is bit-exactly $\pi$-free
(Paper~24~\cite{loutey_paper24}).  \textbf{(H4)} The framework's
results decompose cleanly into a universal algebra-only sector and a
potential-specific Dirac-dependent sector
(Paper~31~\cite{loutey_paper31}).  We discuss the fine-structure
constant combination $K = \pi(B + F - \Delta)$ of
Paper~2~\cite{loutey_paper2_alpha} only in the open-questions section,
in keeping with the project rule that the combination rule remains an
\emph{Observation} at the rule level despite three independent spectral
homes for $B$, $F$, and $\Delta$.

%% ========================================================================
\section{The Packing Construction and the Graph}
\label{sec:packing}
%% ========================================================================

GeoVac's discrete lattice is constructed by a packing argument that is
prior to the Schr\"odinger equation~\cite{loutey_paper0}.  Two axioms
suffice.

\medskip
\noindent\textbf{Axiom 1 (binary distinguishability).}  A minimum of
two distinguishable states is required to define a metric on an
information-bearing manifold.

\medskip
\noindent\textbf{Axiom 2 (geometric indifference).}  In the absence of
external constraints, the distribution of states on a homogeneous
surface maximizes entropy and is therefore isotropic.

\medskip

Applied to concentric isotropic shells in two dimensions, these axioms
force the area-proportional state count on shell~$k$ to be $2k - 1$.
The shell capacities $\{1, 3, 5, 7, \ldots\}$ are isomorphic to the
spherical harmonic basis $Y_{lm}$ under the identification $l = k - 1$
and $m = -l, \ldots, +l$.  The angular quantum numbers $(l, m)$ and
the per-shell degeneracy $2l + 1$ thus emerge from the packing
geometry alone, without invoking the Schr\"odinger equation or the
angular momentum operators $\hat{L}^2$, $\hat{L}_z$.

The principal quantum number~$n$ enters as a cumulative grouping of
the first $n$ shells, producing the count
$\sum_{k=1}^{n} (2k - 1) = n^2$.  Paper~0 is careful to flag that this
grouping is natural but not uniquely forced by the two axioms.
Similarly, the factor of two corresponding to spin multiplicity arises
from the orientability of the compactified packing manifold ($S^2$),
which the paper presents as a structural correspondence rather than a
derivation.  The strong claim is the angular sector $(l, m)$; the
other identifications are framed as structural correspondences.

\subsection*{From the packing to the graph}

The packed lattice acquires a graph structure when the
$SU(2) \otimes SU(1,1)$ ladder operators of the hydrogen dynamical
algebra~\cite{barut1967} are imposed as edges:
\begin{align}
L_\pm |n, l, m\rangle &\propto |n, l, m \pm 1\rangle, \\
T_\pm |n, l, m\rangle &\propto |n \pm 1, l, m\rangle.
\end{align}
The angular operators $L_\pm$ connect states within an $n$-shell along
$m$; the radial operators $T_\pm$ connect states between adjacent
$n$-shells at fixed $l$.  Paper~1~\cite{loutey_paper1} demonstrates
that the spectrum of the resulting graph Hamiltonian, constructed from
the Casimir invariants of the algebra, reproduces the Rydberg formula
$E_n = -1/(2n^2)$ exactly, with no fitting parameters.  At the
algebraic level the discrete and continuous descriptions therefore
already agree on the spectrum.

\subsection*{The universal kinetic constant $\kappa = -1/16$}

A second invariant of the packing geometry is the universal kinetic
constant $\kappa = -1/16$.  Its operational role is to map the
dimensionless $S^3$ eigenvalues $\lambda_n = -(n^2 - 1)$ to the
physical Rydberg spectrum via the energy-shell focal length
$p_0 = Z/n$.  Recent analysis~\cite{loutey_paper18, loutey_paper7} has
clarified its status:\ the value $\kappa = -1/16$ \emph{coincides} with
a geometric quantity computed independently of any energy calibration.
The squared coupling between adjacent $n$-shells in the Gegenbauer
eigenbasis is $c^2(n, l) = (1/16)[1 - l(l+1)/(n(n+1))]$, giving
$c^2(n, 0) = 1/16$ universally for $l = 0$, with equivalent readings
identifying $1/16$ as the squared Chebyshev transition amplitude, as
the inverse Fock Jacobian $1/\Omega^4(0)$, and as the $l = 0$ base rate
of the Casimir decomposition.  That the energy-calibration constant
$\kappa$ \emph{equals} this geometric $1/16$ is recorded as an
\emph{Observation}, not a derivation:\ the geometric value is rigid,
but no mechanism is known that derives $\kappa$'s calibration role from
it (Paper~7~\S III; Paper~18 places $\kappa$ in the \emph{conformal}
tier of the exchange-constant taxonomy accordingly).

\subsection*{Status of the packing-versus-physics question}

The discrete lattice exists \emph{a priori} from the packing
construction, before any reference to the Coulomb potential or the
Schr\"odinger equation.  Its persistent match to known atomic
structure---the Rydberg spectrum at $n_{\max} \to \infty$, the chemical
periodicity of multi-electron atoms (Paper~16, outside the
foundations group), nuclear magic numbers under a harmonic-oscillator
basis (Paper~23, applications group)---is therefore evidence that
quantum-mechanical structure is consistent with a discrete underlying
combinatorial topology, rather than a tautology.  We present this as
the \emph{structural fit} between the packing construction and known
atomic structure.  The synthesis paper does not argue for the
ontological priority of either description; the dual-description
framing of Paper~7 governs the discussion.

%% ========================================================================
\section{The $S^3$ Conformal Equivalence}
\label{sec:s3}
%% ========================================================================

The theoretical foundation of GeoVac is the conformal equivalence
between the discrete Fock graph and the unit $S^3$ Laplace--Beltrami
operator established in Paper~7~\cite{loutey_paper7}.  This section
states the result and explains its operational role.

\subsection*{The Fock projection and the chordal distance identity}

Fock's 1935 construction~\cite{fock1935} writes the momentum-space
hydrogen Schr\"odinger equation
\begin{equation}
\left[ \frac{p^2}{2} - E \right] \psi(\mathbf{p})
- \frac{1}{2\pi^2} \int \frac{\psi(\mathbf{p}')}{|\mathbf{p} - \mathbf{p}'|^2}\,
\mathrm{d}^3 p' = 0
\end{equation}
and applies the stereographic projection
\begin{equation}
u_i = \frac{2 p_0 \, p_i}{p^2 + p_0^2}, \quad
u_4 = \frac{p^2 - p_0^2}{p^2 + p_0^2}, \quad
p_0 = \sqrt{-2 E},
\end{equation}
which maps $\mathbf{p} \in \mathbb{R}^3$ onto the unit three-sphere
$\mathbf{u} \in S^3 \subset \mathbb{R}^4$.  The Coulomb kernel
$1/|\mathbf{p} - \mathbf{p}'|^2$ becomes the inverse-square chordal
distance on $S^3$, and the integral equation reduces to the eigenvalue
problem
\begin{equation}
-\Delta_{S^3}\, \Psi_n = (n^2 - 1)\, \Psi_n,
\end{equation}
free-particle motion on the unit sphere.  The Coulomb potential is
absorbed into the curvature of the sphere; the spectrum is determined
by the topology of $S^3$ alone.

\subsection*{The 18 symbolic proofs}

Paper~7~\cite{loutey_paper7} verifies the conformal-equivalence chain
through 18 independent symbolic proofs implemented in \texttt{sympy}.
The proofs cover the stereographic projection formulas, the unit
sphere constraint, the chordal distance identity, the conformal factor
identity, the conformal Laplacian transformation law, the integer
eigenvalue spectrum, and the energy-shell projection from $S^3$ to
$\mathbb{R}^3$.  These tests live in
\texttt{tests/test\_fock\_projection.py} and
\texttt{tests/test\_fock\_laplacian.py} in the GeoVac repository and
are run as topological integrity checks before any release.  Their
status as machine-verified algebraic identities is a structural
guarantee that the discrete-to-continuum bridge has no hidden steps.

\subsection*{Schr\"odinger recovery in the continuum limit}

The graph Laplacian on the GeoVac lattice converges, in the continuum
limit $n_{\max} \to \infty$ with appropriate rescaling, to the round
$S^3$ Laplace--Beltrami operator.  This convergence, stated empirically
in Paper~7 (verified to $n_{\max} = 30$, without a rate), is now a
theorem:\ for the spinor (Dirac) triple it is an unconditional
van~Suijlekom state-space Gromov--Hausdorff convergence at rate
$(4/\pi)\log n_{\max}/n_{\max}$ under the translation-seminorm
metrization (Paper~38, operator-algebra arc), with the $4/\pi$ constant
universal across compact Lie groups (Paper~40); the scalar
Laplace--Beltrami convergence that Paper~7's 18 proofs target is the
prior-art state-space GH result of
Gaudillot-Estrada--van~Suijlekom (arXiv:2310.14733).  Imposing the energy-shell
constraint $p_0 = Z/n$ then maps the unit-sphere eigenvalues
$\lambda_n = -(n^2 - 1)$ to the hydrogenic Rydberg series
$E_n = -Z^2 / (2 n^2)$.  The dimensionless graph carries no scale; the
Rydberg energy is injected exclusively at the projection step.  This
separation between dimensionless topology and dimensional projection
is the \emph{dimensionless vacuum} viewpoint of Paper~7: the
calibration content of physical observables is concentrated entirely in
the focal-length parameter $p_0$.

\subsection*{Scale-freedom of the substrate and its reflection in conformal field theory}
\label{sec:cft_reflection}

The separation between dimensionless topology and dimensional projection
described above is more than a computational convenience.  It is the
geometric expression of a principle that physics encounters in another
guise:\ conformal field theory.

A conformal field theory is a quantum field theory whose action and
correlators are invariant under rescaling of the underlying coordinates,
together with the larger conformal group $SO(d, 2)$ in $d$ spacetime
dimensions.  Operationally:\ a CFT has no preferred length scale, and
its observables transform predictably under conformal
transformations.  The standard examples include free massless scalars
and fermions on round $S^d$, where the partition functions
$Z_{S^d}^{\text{CFT}}$ are computable in closed form and reproduce
universal coefficients (e.g.\ the F-theorem constant
$F_{S^d}$~\cite{klebanov_pufu_safdi2011}).

GeoVac's substrate exhibits an analogous property at the
\emph{geometric / mathematical} layer, prior to physics.  The packing
construction (Paper~0~\cite{loutey_paper0}) takes a circle of diameter
$d_0$---intended as ``the Planck length'' in spirit, but without
committing to a specific physical value---and outputs a graph whose
Laplacian eigenvalues are the pure rational numbers $n^2 - 1$, the
multiplicities are integers, the Wigner-symbol couplings are exact
rationals, and the packing diameter $d_0$ \emph{cancels out} of every
spectral quantity.  Shrink $d_0$ to the Planck length, grow it to a
galactic radius, set it to anything in between:\ the substrate is the
same.  The eigenvalues that come out of the packing are atomic quantum
numbers in interpretation, but they are not atomic-scale in
construction.  The packing produces a scale-free structural skeleton.

The Fock projection of Paper~7~\cite{loutey_paper7} is a
\emph{conformal map}:\ stereographic projection from $S^3$ to
$\mathbb{R}^3$ preserves angles but rescales lengths by a
position-dependent conformal factor $\Omega(\mathbf{p}) = 2 p_0 /
(p^2 + p_0^2)$.  Physical scales enter the framework through exactly
this factor, in the form of the focal-length parameter $p_0$.  The
scale-free substrate becomes scale-bearing physics only when a specific
$p_0$ is supplied (e.g.\ $p_0 = Z/n$ for hydrogenic atomic levels).
Different choices of $p_0$ correspond to different physical
interpretations of the same combinatorial skeleton; the substrate is
agnostic to the choice.

This is structurally the same principle that CFT scale-invariance
embodies.  A CFT is a physics theory with no preferred length scale at
the level of its action and observables; GeoVac's graph is a
mathematical object with no preferred length scale at the level of its
eigenvalues and matrix elements.  CFT scale-invariance is the physics
side of the principle; the dimensionless graph is the substrate side.
The two converge structurally because they realise the same
property---scale-freedom---at different layers of description, and the
Fock projection that connects them is the conformal map that
implements the scale-free-to-scale-bearing transition.

This reading is not a claim that GeoVac \emph{is} a CFT, nor that the
discrete graph derives CFT scale-invariance.  Those would overshoot:\
GeoVac provides a geometric / spectral substrate, and a CFT is a
quantum field theory living on a substrate.  But the framework's
ability to reproduce continuum CFT-on-sphere partition functions
bit-exactly (recently demonstrated for the conformally coupled scalar
and Camporesi--Higuchi Dirac on $S^3$ and $S^5$ in the operator-algebra
arc, outside the foundations scope of this synthesis but cross-referenced
in the open-questions section below) is structurally explained by the
substrate-side scale-freedom established here.  The framework reproduces
what it does because the substrate is the right kind of mathematical
object to host scale-free physics.

\subsection*{Multi-electron extension to $S^{3N-1}$}

Paper~7~\cite{loutey_paper7} and Paper~21~\cite{loutey_paper21}
extend the construction to $N$-electron atoms by replacing $S^3$ with
the unit sphere $S^{3N-1} \subset \mathbb{R}^{3N}$ under the
$SO(3N)$ rotation group.  The single-electron Fock projection becomes
a multi-electron hyperspherical projection, and the $N$-fermion
antisymmetry condition on $S^{3N-1}$ reads as a representation-theoretic
constraint on the angular sector.  The chemical periodicity of the
periodic table emerges from the $S_N$ subgroup of $SO(3N)$ acting on
the hyperspherical angular basis (see Paper~16, applications group,
outside the foundations scope of this synthesis).

A particular consequence of the $S^3$ projection is the closed-form
single-center two-electron repulsion integral
$F^0(\text{1s}, \text{1s}) = 5/8$, derived from the node-overlap
formula for $V_{ee}$ on $S^3$ in Paper~7.  More generally, all
single-center Slater $F^k$ integrals reduce to exact rationals in
arithmetic over $\mathbb{Q}$.  The single-center two-electron problem
is thus closed in rational arithmetic; transcendentals appear only at
the projection-to-energy step (see Section~\ref{sec:exchange}).

\subsection*{What Paper~7 does and does not claim}

Paper~7 is careful to separate established mathematics from new
contribution.  The stereographic projection formulas, the unit sphere
constraint, the chordal distance identity, the conformal Laplacian
relation, and the integer eigenvalue spectrum are all known results
from Fock, Bargmann, Bander--Itzykson, and Barut--Kleinert.  The new
contributions are: (a) the explicit convergence of the discrete
GeoVac graph to the continuous $S^3$ Laplace--Beltrami operator; (b)
the systematic symbolic verification of the conformal transformation
chain; (c) the framing of physical scales as confined to the
energy-shell projection; (d) the computational consequence that
$O(V)$ sparse graph eigenvalue problems can replace continuous
integration.  Paper~7's explicit statement of the rhetoric---``the
mathematics supports both the discrete-as-fundamental and the
continuous-as-fundamental readings, and we leave the choice to the
reader''---is the foundation of the dual-description framing carried
through the rest of the foundations papers.

%% ========================================================================
\section{Angular Sparsity as a Structural Theorem}
\label{sec:angular}
%% ========================================================================

The angular sparsity theorem of Paper~22~\cite{loutey_paper22} is the
foundations-side statement of the structural mechanism behind GeoVac's
$O(Q^{2.5})$ Pauli-term scaling at the qubit-encoding level.  The
theorem can be stated cleanly without reference to the Coulomb
potential or to any specific radial wavefunction.

\medskip
\noindent\textbf{Theorem (Paper~22).}  \emph{For any $N$-fermion
system whose single-particle orbitals factor as products of spherical
harmonics and arbitrary radial functions, the fraction of nonzero
two-body matrix elements is determined solely by angular momentum
selection rules and is independent of the radial potential.}

\medskip

The proof is a direct application of the Wigner--Eckart theorem and
Condon--Shortley angular momentum
algebra~\cite{CondonShortley1935, Slater1960}.  The two-body matrix
element of a rotationally invariant interaction factors into a sum of
products of Gaunt coefficients (angular factor) and Slater $F^k$,
$G^k$ integrals (radial factor); the Gaunt coefficients vanish for
quartets $(l_a, l_b, l_c, l_d)$ that do not satisfy the
$3j$-triangle inequality
$|l_a - l_c| \le k \le l_a + l_c$ and the corresponding inequality
for $(l_b, l_d)$.  Exact enumeration at fixed $l_{\max}$ yields the
ERI densities reproduced in Table~\ref{tab:eri_density}.

\begin{table}[t]
\centering
\caption{Angular ERI density $D(l_{\max})$ from Paper~22:
fraction of orbital quartets $(a, b, c, d)$ for which the Gaunt
coefficient does not vanish identically.  Universal across Coulomb,
harmonic oscillator, Woods--Saxon, square well, and Yukawa
potentials.}
\label{tab:eri_density}
\begin{tabular}{c c}
\toprule
$l_{\max}$ & $D(l_{\max})$ \\
\midrule
$0$ & $100\%$ \\
$1$ & $7.81\%$ \\
$2$ & $2.76\%$ \\
$3$ & $1.44\%$ \\
$4$ & $0.90\%$ \\
$5$ & $0.62\%$ \\
\bottomrule
\end{tabular}
\end{table}

\subsection*{Universality across potentials}

The bound has been verified numerically across five qualitatively
distinct potentials---Coulomb $-Z/r$, isotropic harmonic oscillator
$\tfrac{1}{2} m \omega^2 r^2$, Woods--Saxon nuclear single-particle
potential, square well, and Yukawa $e^{-\mu r}/r$---and produces
bit-identical sparsity patterns at matched orbital counts.  The
universality is structural: the radial-angular factorization of
$1/r_{12}$ in spherical harmonics (used in quantum chemistry) and the
Talmi--Moshinsky bracket factorization of central NN potentials in
oscillator bases (used in nuclear shell models) impose the same
angular selection rules.  GeoVac's contribution is to state the
explicit quantitative density at each $l_{\max}$ and to frame it as a
universal resource guarantee independent of the radial physics.

\subsection*{Spinor extension}

Paper~22's spinor extension shows that the angular sparsity bound
survives, with mildly inflated prefactor, when the basis is lifted
from $Y_{lm}$ to the spinor spherical harmonics
$\Omega_{\kappa m_j}$ labeled by the Dirac quantum numbers
$(\kappa, m_j)$.  The density $d_\mathrm{spinor}(l_{\max})$ satisfies
$d_\mathrm{spinor} \le d_\mathrm{scalar}$ at every $l_{\max}$ tested,
with the same triangle inequality controlling vanishing of the
matrix elements at the spinor level.  Tier-2 spin-orbit
extensions of the framework therefore inherit the angular sparsity
guarantee without modification.

\subsection*{Operational consequence}

The angular sparsity theorem is what makes GeoVac's qubit-encoded
Hamiltonians competitive for VQE/NISQ resource estimation.  The
quantum-computing applications papers (Group~4) report Pauli-term
counts that scale as $O(Q^{2.5})$ for the composed-geometry
architecture across LiH, BeH$_2$, and H$_2$O---a $51\times$ to
$1712\times$ advantage over published Gaussian baselines at matched
qubit count.  The structural reason for this advantage is the
Gaunt-selection-rule sparsity quantified in Table~\ref{tab:eri_density};
the foundations-paper statement is the theorem itself, not the
specific molecular benchmarks.

%% ========================================================================
\section{Quantum Dynamics on the Graph}
\label{sec:dynamics}
%% ========================================================================

Paper~6~\cite{loutey_paper6} extends the framework from spectroscopy
to time evolution.  The graph Hamiltonian $H$ is sparse with $O(V)$
nonzero entries, so the Crank--Nicolson propagator
\begin{equation}
\psi(t + \Delta t) = \left( I + \tfrac{i}{2} H \Delta t \right)^{-1}
\left( I - \tfrac{i}{2} H \Delta t \right) \psi(t)
\end{equation}
admits implementation at $O(V)$ cost per step using sparse
matrix-vector products plus iterative linear solves.  Unitarity is
preserved to machine precision over $10^4$ time steps in the
$\mathrm{H}_2$ benchmark reported by Paper~6.

The applications demonstrated in Paper~6 are: (a) broadband molecular
spectroscopy from a single delta-kick propagation, recovering 20
dipole-active electronic transitions of H$_2$ with $0.16\%$ mean
error in 33~seconds; (b) resonant Rabi oscillations matching the
analytic Bloch--Siegert-corrected period to $0.41\%$ with $99.98\%$
coherent population transfer; (c) ab initio molecular dynamics with
nuclei coupled to graph forces via Velocity Verlet, conserving total
energy to $0.0003\%$; (d) NVT-ensemble Langevin dynamics with stable
thermal vibrations at 300~K.

Paper~6 is not a foundational claim in the same sense as Papers 7, 18,
22, 24, or 31---it is an operational demonstration that the
foundations support time evolution.  Its role in the synthesis is to
close the gap between the spectral statement of Papers~1 and~7 (the
graph reproduces the Rydberg spectrum) and the dynamical statement
required for any quantum-chemistry application (the graph supports
time evolution at the same $O(V)$ cost).  Without Paper~6's dynamics
engine, GeoVac would be limited to eigenvalue problems; with it, the
framework extends to time-dependent quantum chemistry, spectroscopy,
and molecular dynamics.

%% ========================================================================
\section{The Bargmann--Segal Lattice and the Coulomb/HO Asymmetry}
\label{sec:bs}
%% ========================================================================

Paper~24~\cite{loutey_paper24} extends GeoVac to the three-dimensional
isotropic harmonic oscillator (HO) and uncovers a structural asymmetry
between the Coulomb and HO discretizations.  The construction is built
on the Bargmann--Segal transform~\cite{bargmann1961, hall1994}.

\subsection*{The discrete HO lattice on $S^5$}

The Bargmann--Segal transform maps the 3D isotropic HO Hamiltonian to
the Euler number operator $\hat{N}$ on the Fock--Bargmann space of
holomorphic functions on $\mathbb{C}^3$.  Restriction to the unit
sphere $S^5 \subset \mathbb{C}^3$ identifies the eigenspaces with the
symmetric $SU(3)$ irreducible representations $(N, 0)$, of dimension
$(N + 1)(N + 2)/2$, which match the HO shell degeneracies exactly.
The HO Hamiltonian becomes $\hbar \omega (\hat{N} + 3/2)$, with the
eigenvalues linear in the shell index~$N$.

The discrete graph encoding of this structure takes nodes
$(N, l, m_l)$ to be the labels of the $(N, 0)$-restricted
hyperspherical harmonics on $S^5$, and edges to be dipole transitions
$\Delta N = \pm 1$, $\Delta l = \pm 1$ at exact-rational squared
matrix elements.  At $N_{\max} = 5$ the graph has 56 nodes and 165
edges, all matrix elements computed in exact
\texttt{fractions.Fraction} arithmetic, with zero irrationals
encountered.  The single $\pi$ in the entire Bargmann--Segal
ecosystem is the Gaussian normalization $\pi^{-3}$ of the continuous
Bargmann measure, which the discrete graph never invokes.

\subsection*{The HO rigidity theorem}

\medskip
\noindent\textbf{Theorem (Paper~24, HO rigidity).}  \emph{The 3D
isotropic harmonic oscillator is the unique central potential whose
spectrum arises from the Euler number operator on the Hardy space
$H^2(S^5)$ restricted to symmetric $SU(3)$ irreps $(N, 0)$.}

\medskip

This is the structural dual of the Fock rigidity theorem of Paper~23
(applications group), which states that the $S^3$ stereographic
projection $p_0 = \sqrt{-2 E_n}$ maps a one-electron central-field
Hamiltonian onto the free Laplacian on $S^3$ if and only if the
spectrum is $l$-independent within each $n$-shell, a property unique
to $-Z/r$ by $SO(4)$ symmetry.  Together the two rigidity theorems
establish that the choice of central potential is sharply tied to the
choice of natural manifold: $-Z/r$ to the second-order Riemannian
Laplacian on $S^3$, $\tfrac{1}{2} m \omega^2 r^2$ to the first-order
Euler operator on $H^2(S^5)$.

\subsection*{The Coulomb/HO asymmetry: a six-layer stratification}

A central structural finding of Paper~24, sharpened progressively by
the $S^5$ gauge-structure analysis (Paper~25, QED/gauge group), a
sprint on the wedge KMS state of the framework's Lorentzian extension,
the spectral-action gravity arc (Paper~51), and the master Mellin
engine, is that the Coulomb and HO discretizations differ along
\emph{six} independent layers, identified in stages across 2025--2026.
The asymmetry as currently catalogued reads:

\begin{enumerate}
\item \textbf{Spectrum-computing role of $L_0$.}  On the Coulomb $S^3$
graph, the graph Laplacian $L_0$ \emph{is} the spectrum-generating
operator: $(D - A)$ produces the hydrogenic spectrum $n^2 - 1$.  On
the Bargmann--Segal lattice, $L_0$ is the matter-propagator operator,
not the Hamiltonian: the spectrum comes from the diagonal Euler
operator $\hbar \omega(\hat{N} + 3/2)$, and the graph adjacency
encodes dipole transitions only.

\item \textbf{Calibration $\pi$.}  The Coulomb $S^3$ projection
introduces calibration $\pi$ as the Hopf-base measure (see
Section~\ref{sec:exchange}); the HO Bargmann projection is
bit-exactly $\pi$-free at every finite $N_{\max}$.  This is one of
the clearest concrete instances of an exchange constant being tied
to a specific natural geometry.

\item \textbf{Wilson-physical-content.}  The non-abelian Wilson gauge
construction on $S^3 = SU(2)$ admits natural matter coupling (the
gauge field acts on Pauli spinors carried by the graph nodes); the
parallel SU(3) Wilson construction on the $(N, 0)$ Hardy tower fails
the fixed-group-on-every-link requirement, because transitions
between $(N, 0)$ and $(N + 1, 0)$ irreps are Clebsch--Gordan
intertwiners rather than $SU(3)$ group elements.  The natural gauge
group of the Bargmann graph at the pure-gauge level is $U(1)$, not
$SU(3)$.

\item \textbf{Modular-Hamiltonian structure of the wedge KMS state.}
The Coulomb $S^3$ side admits a hemispheric-wedge KMS state with
closed-form Hilbert--Schmidt orthogonality between the modular
Hamiltonian and the wedge-restricted Dirac operator, with prefactor
$1/\pi^2$ identifying as the Hopf-base measure signature; the HO
Bargmann side does not admit the construction at all.  This is the
operator-algebra-side reading of the asymmetry; for full
details see the operator-algebra arc (Group~1).

\item \textbf{Spectral-action gravity termination.}  The
Camporesi--Higuchi Dirac spectral action on $S^{2m+1}$ terminates at
exactly $(m{+}1)$ power-law terms (Paper~51):\ only $S^3$ ($m = 1$)
gives \emph{pure Einstein gravity} (cosmological $+$ Einstein--Hilbert,
no $R^2$ or higher-curvature corrections), while $S^5$ ($m = 2$)
terminates at three terms including an $R^2$ Gauss--Bonnet
contribution.  The mechanism is the Bernoulli identity for odd
polynomials at half-integer shifts.

\item \textbf{Master Mellin engine chirality-parity.}  On the
Camporesi--Higuchi $S^3$ triple at every finite $n_{\max}$, the Mellin
source $\mathrm{Tr}(D^k e^{-t D^2})$ exhibits a bit-exact
$\mathbb{Z}/2$ chirality-parity factorization (the plain trace
vanishes at odd $k$, the supertrace at even $k$).  On the
Bargmann--Segal Hardy sector this fails on both structural conditions:\
the spectrum is strictly positive (no $\pm$-partner) and no
half-integer-weight chirality grading exists on the $(N, 0)$
tower---propagating the asymmetry to the master-Mellin level.
\end{enumerate}

The first three layers were established in Paper~24 and Paper~25
(QED/gauge group, $S^5$ gauge-structure paper); the fourth (wedge KMS)
was identified in a 2026 sprint and folded into Paper~24's \S{}V; the
fifth (spectral-action gravity) comes from Paper~51; and the sixth
(master Mellin chirality-parity) from a June~2026 sprint.  All six
layers stratify along the same operator-order divide:\ the Coulomb side
uses a second-order Riemannian Laplacian with a nonlinear stereographic
projection, and the HO side uses a first-order complex-analytic Euler
operator with a linear-affine Bargmann projection.

\subsection*{What the asymmetry says about transcendentals}

The bit-exact $\pi$-freeness of the Bargmann--Segal lattice is a
structurally clean statement about where transcendentals enter
physical observables.  $\pi$ is not a universal feature of natural
geometry discretizations; it is a specific structural cost of the
$S^3$ stereographic projection.  This refines the exchange-constant
taxonomy of Paper~18 (see Section~\ref{sec:exchange}): calibration
$\pi$ is Coulomb-specific, tied to second-order Riemannian operators
with nonlinear projections, and has no analog on the first-order
HO side.

%% ========================================================================
\section{The Exchange Constants Taxonomy}
\label{sec:exchange}
%% ========================================================================

Paper~18~\cite{loutey_paper18} catalogues where transcendental
constants enter physical observables computed on the GeoVac graph.
The taxonomy is a six-tier classification organized by the type of
compactness mismatch each transcendental mediates between the discrete
representation and the non-compact coordinate system used for
measurement.

\subsection*{Six tiers}

\begin{enumerate}
\item \textbf{Intrinsic.}  Algebraic content native to the discrete
graph: rational eigenvalues, integer multiplicities, exact-rational
Slater $F^k$ integrals on $S^3$.  $\pi$-free in arithmetic over
$\mathbb{Q}$.

\item \textbf{Conformal / calibration.}  $\pi$ entering through the
Hopf measure of the stereographic projection; the universal kinetic
constant $\kappa = -1/16$; the Hopf-base measure $\text{Vol}(S^2)/4 = \pi$
identifying as the $a_0^2$ Seeley--DeWitt coefficient on the
spinor bundle.  This is the source of calibration $\pi$ in the
Coulomb sector.

\item \textbf{Embedding.}  Transcendentals introduced by embedding the
discrete graph into a continuum coordinate system that the graph does
not natively know: the electron-electron cusp $1/r_{12}$ is the
canonical example.  The cusp is two-dimensional in the angular pair
coordinate $(\alpha, \theta_{12})$ and cannot be reproduced by any
single algebraic insertion; partial-wave Schwartz extrapolation
$\sim -A/(l_{\max} + 2)^4$ is the standard treatment.

\item \textbf{Algebraic-implicit.}  Quantities defined by a polynomial
equation $P(R, \mu) = 0$ with coefficients in a known field extension,
where pointwise diagonalization is a computational convenience rather
than a mathematical necessity.  The Level-3 adiabatic potential curves
$\mu(R)$ on the hyperspherical fiber bundle live in
$\mathbb{Q}(\pi, \sqrt{2})[R]$.

\item \textbf{Composition.}  Single transcendental seeds appearing
through algebraic reduction of integrals: the seed $e^a \cdot E_1(a)$
that closes the Level-2 prolate-spheroidal radial solver via partial
fraction decomposition; the seed $\gamma(1, x)$ closing the Slater
$F^k$ Laguerre-moment decomposition.

\item \textbf{Inner-factor input data.}  Parameter-tied transcendental
content that the framework consumes as Layer-2 input rather than
deriving from Layer-1 graph data.  The Yukawa Dirichlet ring
$\mathbb{Q}[y_1^{-2s}, \ldots, y_n^{-2s}]$ associated with the
inner-factor of a Connes--Marcolli almost-commutative spectral triple
is the canonical example; nuclear-structure scalars ($\langle r^2
\rangle_E$, Zemach radius $r_Z$, quadrupole moment $Q_N$) are also
parameter-tied Layer-2 inputs.  This tier captures the boundary
between what GeoVac derives and what it consumes from external
spectroscopy.
\end{enumerate}

The taxonomy is falsifiable: the claim is that the transcendental
content of any GeoVac observable is determined entirely by the type of
projection linking the compact discrete description to the
non-compact coordinate system.  No observable has been identified that
violates the taxonomy; recent work on bound-state QED in the
hydrogen-Lamb-shift sprint has confirmed that the framework-native
one-loop contributions sit cleanly in the calibration and composition
tiers, with multi-loop UV-renormalization counterterms moving the
relevant content into the inner-factor input data tier.

\subsection*{The master Mellin engine}

A structural sharpening of the taxonomy is the master Mellin engine
identified in Paper~18 \S{}III.7 and formalized in the operator-algebra
arc of the project.  The Weyl--Selberg identification of GeoVac's
exchange constants with Mellin transforms of heat kernels
\begin{equation}
\mathcal{M}\!\left[ \mathrm{Tr}\left( D^k \cdot e^{-t D^2} \right) \right](s)
\end{equation}
indexes three sub-mechanisms by the integer operator order
$k \in \{0, 1, 2\}$:

\begin{itemize}
\item \textbf{M1 ($k = 0$).}  Trivial heat kernel collapses to the
Hopf-base measure $\text{Vol}(S^2)/\pi^2 = 4/\pi$.  This is the
signature of all calibration-$\pi$ content in the Coulomb sector and
appears as the leading constant in the GeoVac propinquity convergence
rate.

\item \textbf{M2 ($k = 2$).}  Standard Seeley--DeWitt expansion of
$\mathrm{Tr}\,e^{-t D^2}$.  Produces $\sqrt{\pi}\,\mathbb{Q} \oplus
\pi^2 \,\mathbb{Q}$ at every integer $s$ via Jacobi $\theta$-inversion
on the half-integer-shifted Dirac spectrum.

\item \textbf{M3 ($k = 1$).}  Mellin transform of
$\mathrm{Tr}(D \cdot e^{-t D^2})$.  Produces the quarter-integer
Hurwitz / vertex-parity content (Catalan's $G$, Dirichlet $\beta(2k)$)
that appears in two-loop QED vertex corrections on $S^3$.
\end{itemize}

The master Mellin engine reading is that M1, M2, M3 are not three
independent mechanisms but three sub-cases of the same Mellin
transform on the GeoVac spectral triple, indexed by which power of the
Dirac operator appears in the trace.  Each sub-mechanism has a
characteristic domain of observable: M1 governs propinquity
convergence rates and Hopf-measure normalizations, M2 governs
heat-kernel and spectral-action coefficients, and M3 governs
vertex-restricted parity-character observables.  Mixing the operator
order of the question with the operator order of the answer produces
the structural degeneracies catalogued in the project's sprint
record.

\subsection*{The Layer-1 / Layer-2 split}

A particularly clean operational consequence of the taxonomy is the
two-layer architecture made explicit by Paper~34 (catalogue paper,
applications group): the bare combinatorial graph (Layer~1, $\pi$-free
in $\mathbb{Q}[\sqrt{d}]$) versus a finite family of named projections
(Layer~2) that add specific variables, dimensions, and transcendentals
when the graph is mapped to a physical observable.  Paper~34
documents 28~projections to date; each appears in the taxonomy of
Paper~18 as either an intrinsic, calibration, embedding,
algebraic-implicit, composition, or inner-factor input data
mechanism.  The two papers are complementary: Paper~18 organizes by
type of compactness mismatch (what \emph{kind} of transcendental
appears), and Paper~34 organizes by named mechanism (\emph{which}
projection appears in \emph{which} observable).

%% ========================================================================
\section{The Universal/Coulomb Partition}
\label{sec:partition}
%% ========================================================================

Paper~31~\cite{loutey_paper31} states the cleanest single
articulation of what GeoVac claims as universal versus what it claims
as potential-tied.  The framework is read as a discrete spectral
triple $(\mathcal{A}, \mathcal{H}, D)$ in the sense of
Connes~\cite{connes_book1994} and the graph-spectral-triple lineage of
Marcolli and van Suijlekom~\cite{marcolli_vs2014, perez_sanchez2024}:
$\mathcal{A}$ is an algebra of functions on the Fock-projected
discrete graph, $\mathcal{H}$ is a Hilbert space of scalar and spinor
states, and $D$ is a Dirac operator inherited from the
Camporesi--Higuchi spectrum on $S^3$~\cite{camporesi_higuchi1996}.

\subsection*{The A/D split}

The structural decomposition reads:

\medskip
\noindent\textbf{Universal sector ($\mathcal{A}$-only).}  Results that
depend only on the algebra $\mathcal{A}$ and the angular structure of
the basis, and \emph{not} on the specific choice of Dirac operator.
These transfer free across central potentials.  Examples:
\begin{itemize}
\item Angular sparsity (Paper~22): ERI density $D(l_{\max})$ depends
only on the spherical-harmonic angular factorization.
\item Composed scaling (Paper~14, applications group): the
$O(Q^{2.5})$ Pauli-term scaling of composed-geometry qubit
encodings depends on the Gaunt selection rules.
\item Selection rules: $\Delta l = \pm 1$, $\Delta m = 0, \pm 1$ for
dipole transitions; the $3j$-triangle inequality at every vertex of
the two-body integral.
\end{itemize}

\medskip
\noindent\textbf{Potential-specific sector ($D$-dependent).}  Results
that depend on the specific spectrum and matrix elements of the Dirac
operator $D$, and therefore on the specific central potential.  These
require per-system rederivation.  Examples:
\begin{itemize}
\item The spectrum itself: $n^2 - 1$ for Coulomb, $N + 3/2$ for HO.
\item The calibration $\kappa = -1/16$ (Coulomb-specific; no analog
on the Bargmann lattice).
\item Transcendental content: $\pi$ for Coulomb via the Hopf measure,
$\pi$-free for HO.
\item Rigidity theorems: Fock rigidity (Paper~23) pins the Coulomb
geometry to $S^3$; HO rigidity (Paper~24) pins the HO geometry to
$H^2(S^5)$ with $(N, 0)$ symmetric irreps.
\item The candidate fine-structure-constant combination
$K = \pi(B + F - \Delta)$ of Paper~2: lives entirely in the
Coulomb-specific sector by its dependence on the Camporesi--Higuchi
Dirac spectrum on $S^3$.
\end{itemize}

\subsection*{Why the partition matters operationally}

The partition organizes Papers 22, 23, 24, and 25 into a single
conceptual structure.  Paper~22's angular sparsity theorem lives in
the universal sector and transfers free between hydrogen and the
deuteron; Paper~23's nuclear shell model qubit Hamiltonians inherit
the angular sparsity automatically but require a fresh per-potential
analysis of the Dirac sector.  Paper~24's HO rigidity and the
$\pi$-free certificate live in the potential-specific sector and
explicitly do \emph{not} transfer to the Coulomb side---the
Bargmann-style first-order operator has no analog on $S^3$.  Paper~25
(synthesis group) on the SU(2) Wilson lattice gauge structure of the
Hopf graph also lives in the potential-specific sector: it depends on
the parallelizability of $S^3 = SU(2)$, which has no analog on $S^5$.

The operational utility of the partition is that anyone porting GeoVac
to a new system can read off, from the universal sector, what
structural properties they inherit for free, and from the
potential-specific sector, what calibrations they need to derive
afresh.  For nuclear shell models, gauge-theoretic extensions, and
oscillator-based quantum dots, the universal sector provides angular
sparsity and composed scaling automatically; the potential-specific
sector has to be redone for each system.

\subsection*{Status as a spectral-triple-language reformulation}

Paper~31 is careful to frame the partition as a structural
reformulation, not a new theorem.  The angular sparsity theorem
(Paper~22), the Fock rigidity theorem (Paper~23), and the HO
rigidity theorem (Paper~24) are each established results; the
partition is the observation that they decompose cleanly along the
algebra-only / Dirac-dependent divide.  The framing uses the language
of NCG and spectral triples because that language has the structural
vocabulary to articulate the divide; the underlying physics content
is unchanged.

%% ========================================================================
\section{The Reconvergence: Two-Body Structure, Periods, and the
Tannakian Substrate}
\label{sec:reconvergence}
%% ========================================================================

The foundations arc described so far \emph{fans out} from a single
taproot.  The packing construction (Paper~0), made physical by the Fock
$S^3$ projection (Paper~7), is the root; from it the framework branches
into spectral graph methods (Paper~1), the exchange-constant taxonomy
(Paper~18), angular sparsity (Paper~22), the Bargmann--Segal lattice
(Paper~24), and the $\mathcal{A}/D$ partition (Paper~31).  Four more
recent papers do the opposite---they \emph{reconverge}, looping the
scattered structural findings back into a single foundational statement
about what the framework autonomously generates and what it consumes.

\subsection{Two-body structure (Paper~54)}
The $\mathcal{A}/D$ partition of Paper~31 is, in the single-particle
setting, a sorting of results into algebra-only and Dirac-dependent.
Paper~54 tests it at the two-body level by constructing the
tensor-product spectral triple $\mathcal{T}_{S^3} \otimes
\mathcal{T}_{S^3}$ and asking whether the gauged spectral action
generates the inter-particle Coulomb interaction from the framework's
own structure.  The free tensor-product Dirac $D_{\mathrm{total}} = D_1
\otimes I + \gamma_1 \otimes D_2$ satisfies $\{D, \gamma\} = 0$ exactly,
so the free spectral action factorizes and contains no
interaction---the correct statement that free particles do not interact.
Gauging via the inner fluctuation of the truncated operator-system
algebra produces a genuine, non-factorizable two-body interaction
(connected fraction $77\%$ at $n_{\max} = 2$).  Its angular structure is
exactly what the universal sector predicts:\ $100\%$ $m$-conserving,
$100\%$ Gaunt-compatible, pure $k = 0$ monopole at both tested cutoffs.
Its radial weights, however, are \emph{not} proportional to the Coulomb
$1/r_{12}$ kernel (Pearson $r = 0.41$--$0.58$, \emph{decreasing} with
$n_{\max}$---structural, not finite-size).  This delineates the
partition at the two-body level:\ the algebra $\mathcal{A}$ determines
\emph{which} channels couple (angular selection rules, universal across
central potentials by Paper~22), and the Dirac operator / metric
determines \emph{how strongly} (radial coupling strengths,
potential-specific).  The recurring ``multi-focal composition wall''
sits exactly on this $\mathcal{A}/D$ boundary, and the reason is
mathematical:\ the spectral action computes Seeley--DeWitt heat-kernel
coefficients (metric curvature), while $1/r_{12}$ is a Green's-function
kernel; the Coulomb potential enters through the Fock projection's
chordal-distance identity (Paper~7), not through the spectral action.

\subsection{Periods (Paper~55)}
The master Mellin engine of Paper~18 indexes the framework's
$\pi$-content by operator order $k \in \{0, 1, 2\}$ in
$\mathcal{M}[\mathrm{Tr}(D^k e^{-t D^2})]$.  Paper~55 completes this
engine into a period-theoretic statement:\ every transcendental that
has appeared in a GeoVac spectral observable---integer powers of $\pi$,
odd zeta values $\zeta(2k{+}1)$, Catalan $G = \beta(2)$ and
Dirichlet-$\beta$ values, and a depth-$k$ loop tower---sits in
cyclotomic mixed-Tate periods at level $\le 4$ over $\mathbb{Z}[i, 1/2]$,
stratified by Mellin slot:\ $M_1 \subset \mathbb{Q}[\pi, \pi^{-1}]$
(pure-Tate, depth~0), $M_2 \subset \bigoplus_k \pi^{2k}\mathbb{Q}$
(pure-Tate even-weight, depth~0, a strict refinement of the
Fathizadeh--Marcolli classification specialized to the static $S^3$
sub-case), and $M_3 \subset \mathcal{MT}(\mathbb{Z}[i, 1/2])$ at level
$\le 4$ (cyclotomic mixed-Tate, via Hurwitz zeta at quarter-integer
shifts and the Deligne--Glanois descent; the vertex-parity-as-$\chi_{-4}$
correspondence of Paper~28).  This is the period-theoretic completion of
the $\mathcal{A}/D$ partition:\ the universal sector's transcendentals
are pinned to named cyclotomic mixed-Tate rings rather than chased
indefinitely.  (Discipline check folded in here:\ the constant
$S_{\min}$, long carried as ``PSLQ-irreducible,'' is now resolved as
$S_{\min} \in \mathbb{Q}[\pi^2, \ln 2, \zeta(3), \zeta(5)]$ in closed
form---the prior label was a basis-coverage artifact.)

\subsection{The Tannakian substrate and the cosmic-Galois injection
(Paper~56)}
If the framework's periods all live in a single named period ring, the
natural categorical home is a motivic Galois group.  Paper~56 assembles,
bit-exact on rational panels (thousands of zero residuals across
$n_{\max} \in \{2, 3, 4\}$), the Tannakian substrate of the GeoVac
spectral triple at finite cutoff:\ a cofiltered pro-system of
finite-cutoff Hopf-algebra truncations, an abelian symmetric-monoidal
rigid representation category, a faithful exact $\otimes$-fiber functor,
and an explicit inclusion of the candidate cosmic-Galois group
$U^*_{\mathrm{Levi}} = \mathbb{G}_a^{3N} \rtimes SL_2$ into
$\mathrm{Aut}^{\otimes}(\omega)$.  The load-bearing result is the
\emph{cosmic-Galois injection}:\ a proven, theorem-grade closed
immersion $U^*_{\mathrm{GV}} \hookrightarrow \mathcal{U}_4^{\mathrm{ab}}
\rtimes SL_2$ of the framework's finite-cutoff Galois group into the
abelianization of the unipotent radical of the level-4 motivic Galois
group $\mathcal{G}_4 = \mathcal{G}_{\mathrm{MT}(\mathbb{Z}[i, 1/2])}$
(Deligne 2010, Glanois 2015).  The comparison target is the
\emph{level-4} group, not Brown's $\mathcal{G}_{\mathrm{MT}(\mathbb{Z})}$:\
Eskandari--Murty--Nemoto 2025 prove Catalan $G$ is not a period of
$\mathrm{MT}(\mathbb{Z})$ but is at level~4, forcing the refinement.
Two scope qualifiers are essential and stated honestly:\ (i) GeoVac
\emph{realizes the abelianization} of the level-4 cosmic Galois group
(Reading~A)---its primitive Hopf substrate is cocommutative, confirmed
bit-exact at the JLO entire-cyclic-cocycle level at depth~2, so the
image lands in $\mathcal{U}_4^{\mathrm{ab}} \cong \mathbb{G}_a^\infty$
rather than the full free non-abelian unipotent; (ii) \emph{infinite-cutoff
equality is not claimed}---the injection direction is closed at theorem
grade at finite cutoff, while the converse equality on depth-$\ge 2$
content is the named multi-year forward question.

\subsection{The forced/free seam (Paper~57)}
The reconvergence closes with an observations paper that consolidates,
across the corpus, the boundary between what the framework forces and
what it consumes.  Paper~57 catalogues the forced entries (graph
spectrum, the $S^3$ outer manifold, the gauge group $U(1) \times SU(2)
\times SU(3)$ from Bertrand $\times$ Hopf-tower truncation, the master
Mellin engine, two-term-exact gravity), the single admitted-not-forced
entry (the quaternionic factor $\mathbb{H}$), and the calibration
entries (Yukawa values, generation count, cutoff moments, atomic
screening exponents, multi-loop QED counterterms).  It identifies the
cross-domain unity pattern---the same multi-focal composition wall
recurs in six independent settings---and tests a packing-reachability
discriminator (an entry is forced iff a witness derivation exists from
the packing axiom plus the real-spectral-triple axioms plus Hopf-tower
truncation plus Bertrand), which reaches $98.3\%$ accuracy across the
$60$-entry catalogue.  The free side splits into two structurally
distinct families:\ a multi-focal composition family (the framework can
build it but cannot autonomously pick a number) and an inner-factor
input-data family (categorical to the outer-factor mechanism; lives
elsewhere by structure).  This is the foundational reconvergence:\ the
$\mathcal{A}/D$ partition (Paper~31), the period classification
(Paper~55), and the Tannakian injection (Paper~56) together turn the
empirical forced/free observation into a theorem-bound boundary.  The
framework's identity is precisely this seam---it maps the structural
skeleton of quantum mechanics (the compact, Peter--Weyl, discrete
regime) and consumes, rather than generates, the continuum calibration
data the skeleton does not reach.

%% ========================================================================
\section{Open Questions and Outlook}
\label{sec:open}
%% ========================================================================

The foundations arc as catalogued here leaves four named structural
questions open.

\subsection*{The $\alpha$ combination rule}

Paper~2~\cite{loutey_paper2_alpha} (observations group) presents the
candidate combination
\begin{equation}
K = \pi (B + F - \Delta), \qquad K \stackrel{?}{=} \alpha^{-1},
\end{equation}
where $B = 42$ is a finite Casimir trace at the truncation $m = 3$,
$F = \zeta_R(2) = \pi^2/6$ is the Fock-degeneracy Dirichlet series at
the packing exponent $d_{\max} = 4$, and $\Delta = 1/40$ is the
inverse Dirac mode count $g_3^{\text{Dirac}}(S^3) = 2(n+1)(n+2)|_{n=3}$.
Each ingredient now has a structurally identified spectral home;
twelve mechanisms have been eliminated as candidate single-principle
derivations of the sum~\cite{loutey_paper18}.  The combination rule
itself remains an \emph{Observation} at the rule level (never a
conjecture or derivation, per the project's hard rule):\ there is no
first-principles derivation explaining why the additive combination of
three structurally independent objects equals $\alpha^{-1}$ at
$8.8 \times 10^{-8}$.  The synthesis treats this as an open structural
question and does not assert priority for the observed identification.

\subsection*{The second packing axiom (W3)}

The packing construction of Paper~0 generates the angular sector
$(l, m)$ and provides natural correspondences for $n$ and spin
multiplicity, but it does not produce calibration data such as Yukawa
couplings or the number of fermion generations.  These appear in the
exchange-constant taxonomy as Layer-2 inner-factor input data: the
framework consumes them rather than deriving them.  The open question
labeled W3 in the project's internal taxonomy is whether a second
packing-axiom-like construction generates calibration data from a
similarly minimal principle.  A 2026 sprint surveying candidates in the
spectral-zeta / master-Mellin-engine ring returned uniformly negative
results across PMNS angles, lepton masses, and a $\sim 21{,}000$-form
mechanical basis~\cite{loutey_paper34}.  The second-packing-axiom
question is open and is structurally distinct from ``find a packing on
a different geometry''---the latter has been answered for the HO via
the Bargmann--Segal lattice without addressing W3.

\subsection*{Cross-manifold extensions}

The Coulomb and HO discretizations live on $S^3$ and the Hardy sector
of $S^5$ respectively.  A natural question is whether a coupled system
(electronic plus nuclear, or two distinct potential regimes within one
molecule) admits a tensor-product spectral triple of the form
$\mathcal{T}_{S^3} \otimes \mathcal{T}_{\mathrm{Hardy}(S^5)}$.  Within
the foundations arc, Paper~24's six-layer Coulomb/HO asymmetry
documents structural obstructions to such a construction at the
NCG-framework level.  The operator-algebra arc (Group~1) has formally
named this as a frontier-of-field open problem; no published framework
currently supports the cross-manifold tensor product.

\subsection*{The multi-focal composition wall}

The framework's composition rules for multiple Fock-style projections
constitute an active research direction.  The internal taxonomy
labels three sub-walls: W1 (within-framework architectural gaps:
cross-register coordinate operators, magnetization-distribution
operators, frozen-core cross-center screening), W2 (multi-loop UV/IR
composition and cross-manifold extensions), and W3 (inner-factor
calibration data, discussed above).  W1a/b/c have been closed at the
first-pass level by the Phase~C work documented in the applications
group; W2a (multi-loop renormalization) remains open at the
field-theory frontier of the operator-algebra arc.  None of these
walls falsify the foundations as presented in this synthesis; they
sharpen the open-question boundary of the framework.

%% ========================================================================
\section{Conclusion}
\label{sec:conclusion}
%% ========================================================================

The Group~3 foundations arc of GeoVac establishes a small set of
structural results that together provide a coherent discretization of
non-relativistic quantum mechanics on natural geometries.  The
discrete Fock graph is provably conformally equivalent to the unit
$S^3$ Laplace--Beltrami operator (Paper~7~\cite{loutey_paper7}); the
angular sparsity of two-body matrix elements is potential-independent
(Paper~22~\cite{loutey_paper22}); the 3D harmonic oscillator admits a
bit-exactly $\pi$-free discrete encoding on $H^2(S^5)$
(Paper~24~\cite{loutey_paper24}); the framework's results decompose
cleanly into an algebra-only universal sector and a Dirac-dependent
potential-specific sector (Paper~31~\cite{loutey_paper31}); the
exchange-constant taxonomy of Paper~18~\cite{loutey_paper18} catalogues
where transcendentals enter physical observables and identifies a
master Mellin engine at three operator orders as the generator of all
known $\pi$-content; the packing axiom of
Paper~0~\cite{loutey_paper0} generates the angular sector geometrically;
the spectral graph theory of Paper~1~\cite{loutey_paper1} and the
$O(V)$ time evolution of Paper~6~\cite{loutey_paper6} provide the
computational machinery.

The framework's concrete advantages are independent of the interpretive
question of which description---discrete or continuous---is more
fundamental.  $O(V)$ sparsity, integer eigenvalues on the unit $S^3$,
angular momentum selection rules baked into the basis, bit-exactly
rational matrix elements on the Bargmann--Segal lattice, and the
universal sparsity bound at fixed $l_{\max}$ are statements about the
discrete graph; their continuum counterparts are theorems about the
Laplace--Beltrami operator or the Euler operator on the corresponding
natural manifold; the conformal equivalence guarantees consistency.
The discrete and continuous readings are dual descriptions connected
by proven maps, and the interpretive question is left open.

The open questions catalogued in Section~\ref{sec:open} are
substantive but do not destabilize the structural results of the
foundations arc.  The $\alpha$ combination rule remains an Observation;
the second-packing-axiom question is open; cross-manifold extensions
are blocked at the NCG-framework level; the multi-focal composition
wall is an active research frontier.  Each of these is a frontier-of-field
question rather than a foundational concern, and the synthesis
treats them as such.

\section*{Acknowledgments}

This synthesis paper consolidates work across the Group~3 foundations
arc of the GeoVac project.  The author thanks the broader open-source
computational quantum-chemistry community for the libraries (NumPy,
SciPy, sympy, mpmath, OpenFermion, Qiskit, PennyLane) that make this
work possible.  The framework's papers and code are publicly archived
on Zenodo with permanent DOIs; this synthesis is a project-archive
document and is not formatted for journal submission.

%================================================================
\begin{thebibliography}{99}
%================================================================

%---- External / classical references ----

\bibitem{fock1935}
V.~Fock,
``Zur Theorie des Wasserstoffatoms,''
Z.\ Phys.\ \textbf{98}, 145 (1935).

\bibitem{bargmann1936}
V.~Bargmann,
``Zur Theorie des Wasserstoffatoms: Bemerkungen zur gleichnamigen Arbeit von V.~Fock,''
Z.\ Phys.\ \textbf{99}, 576 (1936).

\bibitem{bander_itzykson1966}
M.~Bander and C.~Itzykson,
``Group Theory and the Hydrogen Atom,''
Rev.\ Mod.\ Phys.\ \textbf{38}, 330 and 346 (1966).

\bibitem{barut1967}
A.~O.~Barut and H.~Kleinert,
``Transition Probabilities of the Hydrogen Atom from Noncompact Dynamical Groups,''
Phys.\ Rev.\ \textbf{156}, 1541 (1967).

\bibitem{bargmann1961}
V.~Bargmann,
``On a Hilbert space of analytic functions and an associated integral transform, Part I,''
Commun.\ Pure Appl.\ Math.\ \textbf{14}, 187 (1961).

\bibitem{hall1994}
B.~C.~Hall,
``The Segal--Bargmann coherent state transform for compact Lie groups,''
J.\ Funct.\ Anal.\ \textbf{122}, 103 (1994).

\bibitem{camporesi_higuchi1996}
R.~Camporesi and A.~Higuchi,
``On the eigenfunctions of the Dirac operator on spheres and real hyperbolic spaces,''
J.\ Geom.\ Phys.\ \textbf{20}, 1 (1996).

\bibitem{avery_book1989}
J.~Avery,
\textit{Hyperspherical Harmonics: Applications in Quantum Theory}
(Kluwer Academic, Dordrecht, 1989).

\bibitem{aquilanti_caligiana2003}
V.~Aquilanti and A.~Caligiana,
``Sturmian approach to one-electron many-center systems: integrals and iteration schemes,''
Chem.\ Phys.\ Lett.\ \textbf{366}, 157 (2003).

\bibitem{connes_book1994}
A.~Connes,
\textit{Noncommutative Geometry}
(Academic Press, San Diego, 1994).

\bibitem{marcolli_vs2014}
M.~Marcolli and W.~D.~van Suijlekom,
``Gauge networks in noncommutative geometry,''
J.\ Geom.\ Phys.\ \textbf{75}, 71 (2014); arXiv:1301.3480.

\bibitem{perez_sanchez2024}
C.~I.~P\'erez-S\'anchez,
``Bratteli networks and the Spectral Action on quivers,''
arXiv:2401.03705 (2024).

\bibitem{chamseddine_connes1997}
A.~H.~Chamseddine and A.~Connes,
``The Spectral Action Principle,''
Commun.\ Math.\ Phys.\ \textbf{186}, 731 (1997).

\bibitem{CondonShortley1935}
E.~U.~Condon and G.~H.~Shortley,
\textit{The Theory of Atomic Spectra}
(Cambridge University Press, Cambridge, 1935).

\bibitem{Slater1960}
J.~C.~Slater,
\textit{Quantum Theory of Atomic Structure}, Vols.~I and~II
(McGraw-Hill, New York, 1960).

\bibitem{Whitten1973}
J.~L.~Whitten,
``Coulombic potential energy integrals and approximations,''
J.\ Chem.\ Phys.\ \textbf{58}, 4496 (1973).

\bibitem{Dunlap2000}
B.~I.~Dunlap,
``Robust and variational fitting,''
Phys.\ Chem.\ Chem.\ Phys.\ \textbf{2}, 2113 (2000).

\bibitem{Suhonen2007}
J.~Suhonen,
\textit{From Nucleons to Nucleus: Concepts of Microscopic Nuclear Theory}
(Springer, Berlin, 2007).

\bibitem{Peruzzo2014}
A.~Peruzzo et al.,
``A variational eigenvalue solver on a photonic quantum processor,''
Nat.\ Commun.\ \textbf{5}, 4213 (2014).

\bibitem{Lee2021}
J.~Lee et al.,
``Even more efficient quantum computations of chemistry through tensor hypercontraction,''
PRX Quantum \textbf{2}, 030305 (2021).

\bibitem{biedenharn1981}
L.~C.~Biedenharn and J.~D.~Louck,
\textit{Angular Momentum in Quantum Physics}
(Addison-Wesley, Reading, MA, 1981).

%---- GeoVac papers (internal Group 3 series) ----

\bibitem{loutey_paper0}
J.~Loutey,
``Angular Momentum as Information Geometry: Isotropic Packing and the Universal Angular Cross-Section,''
GeoVac Paper~0 (2026).

\bibitem{loutey_paper1}
J.~Loutey,
``The Geometric Atom: Atomic Structure as a Packing Problem,''
GeoVac Paper~1 (2026).

\bibitem{loutey_paper2_alpha}
J.~Loutey,
``The Fine Structure Constant as a Spectral Coincidence on the Hopf Fibration,''
GeoVac Paper~2 (Observations, 2026).

\bibitem{loutey_paper6}
J.~Loutey,
``$O(N)$ Quantum Dynamics and Molecular Spectroscopy via Spectral Graph Theory,''
GeoVac Paper~6 (2026).

\bibitem{loutey_paper7}
J.~Loutey,
``The Dimensionless Vacuum: Recovering the Schr\"odinger Equation from Scale-Invariant Graph Topology,''
GeoVac Paper~7 (2026).

\bibitem{loutey_paper18}
J.~Loutey,
``Spectral-Geometric Exchange Constants: Projecting Discrete Spectra onto Continuous Manifolds,''
GeoVac Paper~18 (2026).

\bibitem{loutey_paper21}
J.~Loutey,
``The Geometric Vacuum: Quantum Structure from Compact Manifold Topology,''
GeoVac Paper~21 (Synthesis, 2026).

\bibitem{loutey_paper22}
J.~Loutey,
``Potential-Independent Angular Sparsity for Qubit Hamiltonians in Spherical Fermion Systems,''
GeoVac Paper~22 (2026).

\bibitem{loutey_paper24}
J.~Loutey,
``The Bargmann--Segal Lattice: $\pi$-Free Discretization of the Harmonic Oscillator on $S^5$,''
GeoVac Paper~24 (2026).

\bibitem{loutey_paper31}
J.~Loutey,
``The Universal/Coulomb Partition of the GeoVac Spectral Triple,''
GeoVac Paper~31 (2026).

\bibitem{loutey_paper34}
J.~Loutey,
``GeoVac as a Two-Layer Framework: Projection Taxonomy and the Empirical Matches Catalog,''
GeoVac Paper~34 (Observations, 2026).

\bibitem{klebanov_pufu_safdi2011}
I.~R.~Klebanov, S.~S.~Pufu, B.~R.~Safdi,
``F-Theorem without Supersymmetry,''
\textit{Journal of High Energy Physics} \textbf{10} (2011) 038,
arXiv:1105.4598.

\end{thebibliography}

\end{document}
